\section{LifecycleGate Semantic Benchmarks}
\label{sec:appendix-benchmarks}

This appendix contains the low-dimensional semantic benchmarks that
distinguish current-action filtering from update admission.  They instantiate
the certificate semantics of Section~\ref{sec:gate} rather than the
authority-transition attack evaluated in Section~\ref{sec:study}.

\subsection{Setup}

We instantiated the linear system of Appendix~\ref{sec:appendix-lift} with
$\lambda=1.2$,
$u_{\max}=0.2$, and $S=[-1,1]$.  FDI induced the ambiguity set
$X_{\mathcal A}(h)=[-\rho,\rho]$.  The common certified kernel for invariant
certificate $I=S$ is
\begin{equation}
    K(\rho)=[-u_{\max},u_{\max}]\cap[-1+\lambda\rho,1-\lambda\rho].
\end{equation}
The kernel is nonempty exactly when $\lambda\rho\le 1$.

\subsection{Learning-Freeze Frontier}

The frontier occurs at $\rho=1/\lambda=0.833$. Table~\ref{tab:frontier} shows the exact kernel values for representative ambiguity radii.

\begin{table}[t]
\centering
\caption{Certified kernel for the one-dimensional ambiguity set.}
\label{tab:frontier}
\begin{tabular}{cccc}
\toprule
$\rho$ & $\lambda\rho$ & $K(\rho)$ & Gate \\
\midrule
0.600 & 0.720 & $[-0.200,0.200]$ & allow \\
0.800 & 0.960 & $[-0.040,0.040]$ & constrain \\
0.833 & 1.000 & $\{0\}$ & boundary \\
0.900 & 1.080 & $\emptyset$ & freeze/reject \\
1.000 & 1.200 & $\emptyset$ & freeze/reject \\
\bottomrule
\end{tabular}
\end{table}

At $\rho=0.9$, states $+0.9$ and $-0.9$ are individually controllable when
their sign is trusted.  Their safe intervals are $[-0.2,-0.08]$ and
$[0.08,0.2]$, respectively.  The empty intersection prevents a common
history-dependent action from being certified.

\subsection{Trusted Anchors and Learning-Aware FDI}

A trusted state anchor can restore learning by shrinking the plausible-state
set.  At $\rho=0.9$, a trusted sign anchor reduces the set to $[0,0.9]$ or
$[-0.9,0]$.  The corresponding kernels, $[-0.2,-0.08]$ and $[0.08,0.2]$,
are nonempty.

Learning-aware FDI creates the complementary failure mode.  For a candidate
action $u=0.1$, the gate allows the update at $\rho=0.6$, projects it to
$0.04$ at $\rho=0.8$, and rejects it at $\rho=0.9$.  The attacked history
therefore determines whether the same candidate is allowed, constrained, or
rejected.

\subsection{Lifecycle Benchmark Results}

We compared five mechanisms over five update steps: an ungated learner, a
nominal action filter, a robust action filter, always-freeze, and LifecycleGate
with projection.  The one-dimensional benchmark used the system above.  The
linear tank benchmark used the level-deviation model
\begin{equation}
    x_{t+1}=Ax_t+Bu_t,
    \quad
    A=\begin{bmatrix}1.08&0.04\\0.02&1.03\end{bmatrix},
    \quad
    B=\begin{bmatrix}1\\0.1\end{bmatrix},
\end{equation}
with $|u_t|\le 0.15$ and $x_t\in[-1,1]^2$. The nonlinear tank benchmark uses square-root inter-tank and drain flows,
\begin{equation}
\begin{aligned}
h_{1,t+1} &= h_{1,t}+\Delta t(q_0+u_t-q_{12}(h_t)-d_1(h_{1,t})),\\
h_{2,t+1} &= h_{2,t}+\Delta t(q_{12}(h_t)-d_2(h_{2,t})),
\end{aligned}
\end{equation}
where $q_{12}$ and $d_i$ are square-root flow terms.  Its interval observer
returned a box around the nominal attacked history.  The one-step kernel was
computed over a grid within that box.  In all benchmarks, learning-aware FDI
pushed the learner toward a larger positive action parameter.

\begin{figure*}[t]
\centering
\includegraphics[width=\textwidth]{lifecycle_gate_summary.pdf}
\caption{Lifecycle-gate benchmark results.  Increasing ambiguity shrinks the
certified kernels of the one-dimensional, linear-tank, and nonlinear-tank
models.  LifecycleGate constrains or freezes updates when action-only filters
would leave the updated policy outside the kernel.}
\label{fig:lifecycle-benchmarks}
\end{figure*}

Figure~\ref{fig:lifecycle-benchmarks} separates current-action filtering from
update certification.  At $\rho=0.8$, the nominal filter had unsafe certified
rate $1.0$.  The robust action filter had stale-policy rate $1.0$ because the
learner parameter left $K(\rho)$ after action projection.  LifecycleGate had
zero unsafe or stale operation and retained one of five updates.  At
$\rho=0.9$, the kernel was empty and LifecycleGate froze learning.

The linear tank benchmark reproduced the same pattern.  At $\rho=0.8$,
nominal filtering had unsafe certified rate $0.8$, and robust action filtering
had stale-policy rate $0.8$.  LifecycleGate kept both rates at zero and allowed
two of five updates.  It allowed one constrained update at $\rho=0.85$ and
froze at the empty-kernel point $\rho=0.92$.

The nonlinear tank benchmark preserved the separation with kernels computed
from sampled nonlinear dynamics.  At $\rho=0.2$, nominal filtering was unsafe
in four of five steps, and robust filtering was stale in four.  LifecycleGate
kept both rates at zero while allowing two constrained updates.  At
$\rho=0.45$, the interval-observer box admitted no common safe action and
LifecycleGate froze the update.

\subsection{Parameter-Level Update Gate}

We also instantiated the gate at the parameter level for linear policy
$u=kz+b$.  Each future attacked observation $z$ induced a safe-action kernel
over ambiguity interval $[z-\rho,z+\rho]$.  Requiring $kz+b$ to lie in every
kernel produced linear halfspace constraints on $(k,b)$.  The parameter gate
projected poisoned candidates onto this certified set.

\begin{figure}[t]
\centering
\includegraphics[width=\columnwidth]{parameter_gate_summary.pdf}
\caption{Parameter-level gate for linear policy $u=kz+b$.  Current-action
filtering leaves the future action map stale, whereas projection onto the
certified halfspace set constrains the full update.}
\label{fig:parameter-gate}
\end{figure}

Figure~\ref{fig:parameter-gate} shows the parameter-level frontier.  At
$\rho=0.7$, current-action filtering had stale-policy rate $1.0$.  The
parameter gate kept the rate at zero and permitted two of five updates.  It
permitted one constrained update at $\rho=0.8$ and froze at $\rho=0.85$.
Figure~\ref{fig:matrix-parameter-gate} repeats the comparison for the matrix
parameterization.

\begin{figure}[t]
\centering
\includegraphics[width=\columnwidth]{matrix_parameter_gate_summary.pdf}
\caption{Parameter-level gate for the two-dimensional linear tank policy.  The
certified set intersects future-history halfspaces, and the gate projects the
poisoned update into that set.}
\label{fig:matrix-parameter-gate}
\end{figure}

We applied the same construction to policy $u=k_1z_1+k_2z_2+b$.  Each grid
observation induced a plausible-state box and a one-step safe-action interval.
These intervals produced halfspaces in $(k_1,k_2,b)$.  At $\rho=0.5$,
current-action filtering had stale-policy rate $1.0$.  The parameter gate kept
the rate at zero and admitted all five projected updates.

\subsection{Adaptive Attacker Stress Test}

We next allow a finite-search white-box attacker to select among poisoned
update directions.  Its objectives maximize stale-policy, unsafe-action, or
forced-freeze outcomes.  For forced freeze, the attacker also selects the
budget-feasible ambiguity radius that most reduces the certified set.

\begin{figure}[t]
\centering
\includegraphics[width=0.98\columnwidth]{adaptive_attack_summary.pdf}
\caption{Finite-search attacker against the scalar parameter gate.  The
attacker selects update directions or ambiguity radii for three objectives.
LifecycleGate freezes when the certified kernel is empty.}
\label{fig:adaptive-attacker}
\end{figure}

Figure~\ref{fig:adaptive-attacker} preserves the lifecycle separation under
adaptive direction search.  At high ambiguity, the robust action filter
protected the applied action but left the raw policy stale.  LifecycleGate
kept unsafe, stale, and uncertified-learning rates at zero while allowing
feasible projected updates.  At $\rho=0.84$, the forced-freeze objective made
the common kernel empty.  The robust filter recorded five uncertified learner
updates, whereas LifecycleGate performed none.

\subsection{Safe-Control-Gym Cartpole Semantics}

We next instantiate the lifecycle condition in the Safe-Control-Gym cartpole
benchmark.  The official linear-quadratic regulator (LQR) is augmented with a
two-parameter linear residual learner.  The attacker injects continuous
observation FDI and optimizes a cross-entropy-method (CEM) objective over
stale-policy, unsafe-action, and forced-freeze margins.  The attacked
observation $z$ and radius $\rho$ induce the plausible box
$X_{\mathcal A}(z)=\{x:\|x-z\|_\infty\le \rho\}$.  We sample this box and
evaluate the official CasADi cartpole dynamics over an action grid to compute
$K_\Gamma(z)$.  A reference run uses the official cartpole CBF and its sampled
Lie-derivative condition.

This experiment is a direct semantic implementation of
Theorem~\ref{thm:gate} and Proposition~\ref{prop:parameter-gate}. The
future-history abstraction $\mathcal Z$ is a local grid of future attacked
observations around the current attacked observation. For each
$z'\in\mathcal Z$, the robust kernel produces an interval
$K_\Gamma(z')=[\ell_{z'},u_{z'}]$. The LQR term is fixed, while the residual
learner is linear in $(\theta_{\mathrm{gain}},\theta_{\mathrm{bias}})$, so
requiring the updated residual policy to stay inside all future kernels yields
the halfspace constraints used by LifecycleGate.

\begin{figure*}[t]
\centering
\includegraphics[width=\textwidth]{safe_control_gym_paper_mini_summary.pdf}
\caption{Safe-Control-Gym cartpole lifecycle evidence from a two-run sweep.
The panels compare the CasADi robust kernel, an Euler attack-search surrogate,
and a sampled CBF reference.  At $\rho=0.03$, LifecycleGate eliminated stale
and uncertified updates while retaining update availability.  The kernel became
empty as ambiguity increased.}
\label{fig:safe-control-gym}
\end{figure*}

Figure~\ref{fig:safe-control-gym} separates three effects.  Action filtering
removed current constraint violations without certifying the updated residual
policy.  Always-freeze avoided stale certification but admitted no learning.
LifecycleGate projected feasible poisoned updates into the future-history
halfspace set, preserving update availability before the freeze frontier.  At
larger ambiguity radii, the kernel became empty and LifecycleGate froze the
update.  These results instantiate the certificate semantics in a standard
control benchmark; the evaluation does not represent an industrial deployment.
